% =====================================================================
%  TYPE: Postoperative instructions -- cervical (ACDF, revision, foraminotomy)
%  Sections: patient block, instructions, follow-up, emergency
%  box, when to call, medication page. Based on Mantarakis 2026-08 and Giannopoulos 2026-06.
%  Written for the anterior approach. Posterior approach: delete the
%  'Κατάποση -- Φωνή' item and change 'πρόσθια επιφάνεια του τραχήλου'.
%  Dural tear / CSF leak: enable the three blocks marked DURAL TEAR.
%
%  Copy, rename to "YYYYMMDD - surname - topic.tex", fill in and
%  compile with XeLaTeX. Toggle alternatives by moving the "%".
%  Text in [brackets] is a placeholder: replace or delete it.
% =====================================================================
\documentclass[10pt,greek]{article}
\usepackage[useregional]{datetime2}

%%%%%%%%%%%%%%%%%%%%% Language
\providecommand\DocLang{gr}

%%%%%%%%%%%%%%%%%%%%% Clinic (letterhead logo, footer address)
\providecommand\Clinic{Peristeri}
% \providecommand\Clinic{Mandra}
% \providecommand\Clinic{Bioclinic}
% \providecommand\Clinic{Psychikou}
% \providecommand\Clinic{Metropolitan}
% \providecommand\Clinic{Dunant}
% \providecommand\Clinic{Euroclinic}

%%%%%%%%%%%%%%%%%%%%% Document code for the QR: paste the code issued
%%%%%%%%%%%%%%%%%%%%% by doc.erns.eu (empty = no QR code)
\providecommand\DocCode{}

%%%%%%%%%%%%%%%%%%%%% Patient
\providecommand\PatientPrefix{Ο}
% \providecommand\PatientPrefix{Η κα}
\providecommand\PatientName{Παπαδόπουλος Γιώργος}
\providecommand\PatientAge{45}
% \providecommand\PatientAge{}

%%%%%%%%%%%%%%%%%%%%% Date
\providecommand\DocDate{\DTMdate{2026-01-15}}
% \providecommand\DocDate{\DTMdate{2026-05-04}}

%%%%%%%%%%%%%%%%%%%%% Title
\providecommand\DocTitle{Μετεγχειρητικές Οδηγίες}
% \providecommand\DocTitle{Ενημερωτικό Σημείωμα για τον Ασθενή}

%%%%%%%%%%%%%%%%%%%%% Closing (empty = no signature block)
\providecommand\DocClosing{Ο ιατρός}
% \providecommand\DocClosing{}

%%%%%%%%%%%%%%%%%%%%% Stamp and signature image
\providecommand\DocStamp{no}
% \providecommand\DocStamp{yes}

\input{evrtemplate}

\begin{document}

\begin{PatientInfo}
  \Info{Διάγνωση}{Τιμή 1}
  \Info{Χειρουργείο}{(10/01/2026) Τιμή 2}
%   \Info{Προηγηθέν}{Τιμή 4}
%   \Info{Συνοδά}{Τιμή 5}
\end{PatientInfo}

Ακολουθούν οδηγίες για την ομαλή σας ανάρρωση.
% DURAL TEAR: replace the line above with:
% Ακολουθούν οδηγίες για την ομαλή σας ανάρρωση. \textbf{Οι οδηγίες 1 και 2 είναι οι πλέον κρίσιμες} και αφορούν την προστασία της σύγκλεισης της μήνιγγας.

% ===== DURAL TEAR (1/3): critical instructions =========================
% \section*{Κρίσιμες Οδηγίες}
% \begin{enumerate}[label=\arabic*.]
% \setlength{\parskip}{6pt}
% \item \textbf{Θέση σώματος -- Αποφυγή οριζόντιας κατάκλισης}
%        \begin{itemize}
% \item \textbf{Αποφύγετε την οριζόντια (ξαπλωτή) θέση όσο είναι δυνατόν.}
% \item Ο ύπνος πρέπει να γίνεται σε \textbf{ημικαθιστή θέση}, με ανύψωση της κεφαλής και του κορμού κατά \textbf{30--45 μοίρες} (π.χ.\ με ανασηκωμένο το κρεβάτι, με πολλαπλά μαξιλάρια στην πλάτη ή σε αναπαυτική πολυθρόνα).
% \item Κατά τη διάρκεια της ημέρας, προτιμάτε την καθιστή ή όρθια θέση.
% \item Η θέση αυτή μειώνει την πίεση του εγκεφαλονωτιαίου υγρού στο σημείο της σύγκλεισης και προστατεύει τη ραφή της μήνιγγας.
% \item Η οδηγία αυτή ισχύει για \textbf{τουλάχιστον 2 εβδομάδες} ή μέχρι νεότερη ιατρική οδηγία.
%     \end{itemize}
% \item \textbf{Αποφυγή αύξησης της ενδοκράνιας πίεσης}
%
% Οποιαδήποτε προσπάθεια αυξάνει απότομα την πίεση του εγκεφαλονωτιαίου υγρού και μπορεί να διαρρήξει τη σύγκλειση. Για τουλάχιστον \textbf{6 εβδομάδες}:
%        \begin{itemize}
% \item \textbf{Καμία άρση βάρους} -- απολύτως τίποτε βαρύτερο από \textbf{2 κιλά} (ούτε σακούλες, ούτε βαλίτσες, ούτε παιδιά, ούτε έπιπλα).
% \item \textbf{Μην σπρώχνετε και μην τραβάτε} αντικείμενα (πόρτες, καρέκλες, καρότσι σούπερ μάρκετ).
% \item \textbf{Μην ζορίζεστε στην τουαλέτα.} Χρησιμοποιήστε το καθαρτικό που σας χορηγήθηκε ώστε οι κενώσεις να είναι μαλακές και χωρίς προσπάθεια.
% \item Εάν χρειαστεί να \textbf{βήξετε ή να φτερνιστείτε}, κάντε το με \textbf{ανοιχτό το στόμα}. Ποτέ με κλειστό στόμα ή κλειστή μύτη.
% \item \textbf{Μην φυσάτε τη μύτη σας.} Σκουπίστε την απαλά.
% \item Αποφύγετε το σκύψιμο με το κεφάλι χαμηλά (π.χ.\ για να δέσετε παπούτσια, να μαζέψετε κάτι από το πάτωμα) -- καθίστε ή ζητήστε βοήθεια.
% \item Καμία έντονη σωματική άσκηση, τρέξιμο, ή αθλητική δραστηριότητα.
%     \end{itemize}
% \end{enumerate}
% =======================================================================

\section*{Γενικές Οδηγίες}
\begin{enumerate}[label=\arabic*.]
\setlength{\parskip}{6pt}
  \item \textbf{Κινητοποίηση}
       \begin{itemize}
\item Επιτρέπεται και συνιστάται η βάδιση εντός της οικίας, με ήπιο ρυθμό και συχνά μικρά διαλείμματα.
\item Κατά την έγερση από το κρεβάτι, γυρίστε πρώτα στο πλάι και στη συνέχεια ανασηκωθείτε \textbf{με τα χέρια}, χωρίς απότομες κινήσεις και χωρίς κοιλιακή προσπάθεια.
\item Χρησιμοποιείτε το αυχενικό κολάρο σύμφωνα με τις οδηγίες του θεράποντος ιατρού.
\item Αποφύγετε για τουλάχιστον 6 εβδομάδες: έντονες κάμψεις, εκτάσεις και στροφές του αυχένα.
\item Αποφύγετε την παρατεταμένη κάμψη του αυχένα προς τα κάτω (π.χ.\ κινητό, ανάγνωση, εργασία σε γραφείο).
\item Αποφύγετε τις απότομες κινήσεις του κεφαλιού.
    \end{itemize}
  \item \textbf{Πόνος -- Φάρμακα}
      \begin{itemize}
\item Λαμβάνετε τα παυσίπονα όπως έχουν συνταγογραφηθεί.
\item Επιδείνωση ή αλλαγή του χαρακτήρα του πόνου πρέπει να αναφέρεται άμεσα.
% \item \textbf{Μην λαμβάνετε αντιφλεγμονώδη} (π.χ.\ Voltaren, Ponstan, Naprosyn, Brufen) χωρίς ιατρική οδηγία.
    \end{itemize}
  \item \textbf{Κατάποση -- Φωνή}
      \begin{itemize}
\item Είναι αναμενόμενη ήπια δυσκολία στην κατάποση και αίσθημα «ξένου σώματος» στον λαιμό τις πρώτες ημέρες έως εβδομάδες. Υποχωρεί σταδιακά.
\item Ξεκινήστε με μαλακή τροφή και υγρά· προχωρήστε σταδιακά σε κανονική δίαιτα ανάλογα με την ανοχή σας.
\item Τρώτε αργά, σε μικρές μπουκιές, σε καθιστή θέση.
\item Πιθανό ήπιο βράχνιασμα ή αλλοίωση της φωνής. Αποφύγετε την υπερβολική χρήση της φωνής (φώναγμα, παρατεταμένη ομιλία) τις πρώτες εβδομάδες.
\item \textbf{Αναφέρετε άμεσα:} σοβαρή ή επιδεινούμενη δυσκολία κατάποσης, πνιγμονή, ή βράχνιασμα που δεν βελτιώνεται μετά τις 2 εβδομάδες.
    \end{itemize}
  \item \textbf{Φροντίδα Τραύματος}
      \begin{itemize}
\item Το τραύμα (πρόσθια επιφάνεια του τραχήλου) πρέπει να παραμένει καθαρό και στεγνό.
\item Αλλαγή τραύματος κάθε 2 ημέρες με αυτοκόλλητα γαζάκια για 7 ημέρες.
\item Αποφύγετε διαβροχή της περιοχής όσο υπάρχουν γάζες.
\item Μετά τις 7 ημέρες μπορεί να χρησιμοποιείται spray τύπου Pulvo47 ή άλλο αντίστοιχο, 2 φορές την ημέρα.
\item Παρακολουθείτε καθημερινά την περιοχή για επανεμφάνιση οιδήματος (πρηξίματος) ή για εκροή υγρού.
    \end{itemize}
% ===== DURAL TEAR (2/3): headache and lumbar drain site ================
%   \item \textbf{Κεφαλαλγία}
%       \begin{itemize}
% \item Ήπια κεφαλαλγία μετά την αφαίρεση της οσφυϊκής παροχέτευσης είναι αναμενόμενη και υποχωρεί σταδιακά εντός ημερών.
% \item Βοηθούν η επαρκής ενυδάτωση, η ξεκούραση και τα απλά παυσίπονα.
% \item \textbf{Αναφέρετε άμεσα} κεφαλαλγία που εμφανίζεται ή επιδεινώνεται σαφώς όταν καθίσετε ή σταθείτε όρθιος και ανακουφίζεται όταν ξαπλώσετε, ιδίως αν συνοδεύεται από ναυτία ή εμβοές.
%     \end{itemize}
%   \item \textbf{Φροντίδα σημείου οσφυϊκής παροχέτευσης (πλάτη)}
%       \begin{itemize}
% \item Το σημείο εισόδου της παροχέτευσης στην οσφυϊκή χώρα παραμένει καλυμμένο με αποστειρωμένο αυτοκόλλητο για 3--5 ημέρες.
% \item Αλλαγή του αυτοκόλλητου εάν λερωθεί ή βραχεί.
% \item \textbf{Αναφέρετε άμεσα} εκροή διαυγούς υγρού, ερυθρότητα, οίδημα ή πόνο στο σημείο αυτό.
%     \end{itemize}
% =======================================================================
  \item \textbf{Διατροφή -- Λειτουργία εντέρου}
    \begin{itemize}
\item Διατροφή πλούσια σε πρωτεΐνες και φυτικές ίνες, μαλακής σύστασης το πρώτο διάστημα.
\item Επαρκής ενυδάτωση -- τουλάχιστον 1,5--2 λίτρα υγρών ημερησίως.
% \item \textbf{Λαμβάνετε καθημερινά το καθαρτικό} που σας χορηγήθηκε. Στόχος είναι κενώσεις μαλακές, χωρίς καμία προσπάθεια.
\item Ενημερώστε τον ιατρό εάν παραμένει δυσκοιλιότητα.
    \end{itemize}
  \item \textbf{Οδήγηση}

Αποφύγετε την οδήγηση για τουλάχιστον 4 εβδομάδες ή μέχρι νεότερη ιατρική οδηγία. Η οδήγηση απαιτεί πλήρες εύρος κινήσεων του αυχένα για ασφαλή έλεγχο του περιβάλλοντος χώρου. Μην οδηγείτε εφόσον:
    \begin{itemize}
\item χρησιμοποιείτε αυχενικό κολάρο,
\item λαμβάνετε φάρμακα που προκαλούν υπνηλία,
\item έχετε σημαντικό πόνο, δυσκαμψία στον αυχένα ή κεφαλαλγία.
    \end{itemize}
  \item \textbf{Ταξίδια}
    \begin{itemize}
\item Αποφύγετε τα αεροπορικά ταξίδια για τουλάχιστον 4--6 εβδομάδες ή μέχρι νεότερη ιατρική οδηγία.
\item Αποφύγετε τα μακρινά οδικά ταξίδια το πρώτο διάστημα.
    \end{itemize}
\end{enumerate}

\section*{Εξειδικευμένες Οδηγίες}
\begin{enumerate}[label=\arabic*.]
  \item \textbf{Αντιβίωση}
      \begin{itemize}
\item \textbf{Zinadol 500 mg}, 1 δισκίο 2 φορές την ημέρα, για \textbf{5 ημέρες}.
\item Ολοκληρώστε το σχήμα ακόμη και αν αισθάνεστε καλά.
    \end{itemize}
  \item \textbf{Αντιθρομβωτική αγωγή}
      \begin{itemize}
\item Κάθε ημέρα Τιμή 6 υποδορίως, για Τιμή 7 ημέρες, την ίδια ώρα.
\item Αποφύγετε μάλαξη στο σημείο ένεσης.
% \item Η τακτική ήπια βάδιση συμβάλλει στην πρόληψη θρόμβωσης.
    \end{itemize}
  \item \textbf{Πώρωση της σπονδυλοδεσίας}
      \begin{itemize}
\item Η οστική πώρωση της σπονδυλοδεσίας απαιτεί χρόνο και ιδιαίτερη προσοχή.
% \item Πρόκειται για \emph{αναθεώρηση} σπονδυλοδεσίας. Η οστική πώρωση απαιτεί χρόνο και ιδιαίτερη προσοχή.
\item \textbf{Απόλυτη διακοπή του καπνίσματος} -- η νικοτίνη μειώνει σημαντικά την πιθανότητα πώρωσης.
% \item Αποφύγετε τη λήψη αντιφλεγμονωδών (ΜΣΑΦ).
    \end{itemize}
  \item \textbf{Φυσικοθεραπεία}
      \begin{itemize}
\item \textbf{Δεν επιτρέπεται} φυσικοθεραπεία, μάλαξη ή κινησιοθεραπεία αυχένα το πρώτο διάστημα.
\item Τον πρώτο μήνα συνιστάται μόνο ήπιο περπάτημα και ελαφριά καθημερινή δραστηριότητα.
\item Ασκήσεις κινητοποίησης ή ενδυνάμωσης θα αποφασιστούν κατά την επανεξέταση.
    \end{itemize}
    \item \textbf{Μπάνιο -- Υγιεινή}
        \begin{itemize}
\item Μπάνιο επιτρέπεται με χρήση αδιάβροχων αυτοκόλλητων γαζών. Αποφύγετε τη διαβροχή του τραύματος. Εάν αυτό συμβεί, στεγνώστε καλά με αποστειρωμένες γάζες και καλύψτε με αυτοκόλλητο.
\item Ολοκληρωτικό ντους επιτρέπεται μετά από 10 ημέρες και μετά από σύσταση του ιατρού.
% \item \textbf{Απαγορεύεται} το μπάνιο σε μπανιέρα, πισίνα ή θάλασσα.
        \end{itemize}
\end{enumerate}

\section*{Επανέλεγχοι}
\begin{description}
\item[Σε 5--7 ημέρες:] κλινική εξέταση και επισκόπηση της τραχηλικής χώρας.
\item[Σε 1 μήνα:] επανεκτίμηση, απεικονιστικός έλεγχος και οδηγίες για δραστηριότητες / φυσικοθεραπεία.
\end{description}

\begin{UrgentBox}
% ===== DURAL TEAR (3/3): CSF leak signs ================================
% \item[] \textbf{Α. Σημεία διαφυγής εγκεφαλονωτιαίου υγρού (ΕΝΥ):}
% \item \textbf{Εκροή διαυγούς, άχρωμου υγρού} από το τραύμα του τραχήλου ή από το σημείο της παροχέτευσης στην πλάτη.
% \item \textbf{Επανεμφάνιση ή αύξηση του οιδήματος} (πρηξίματος) στην τραχηλική χώρα, ή αίσθημα «κύματος» υγρού κάτω από το δέρμα.
% \item \textbf{Έντονη κεφαλαλγία που επιδεινώνεται στην όρθια θέση} και ανακουφίζεται στην κατάκλιση.
% \item Αλμυρή ή μεταλλική γεύση στο στόμα, ή εκροή διαυγούς υγρού από τη μύτη.
% =======================================================================
\item[] \textbf{Σημεία λοίμωξης / μηνιγγίτιδας:}
\item \textbf{Πυρετός} (Θερμοκρασία $>$ 38\textdegree{}C) με \textbf{έντονη κεφαλαλγία}.
\item \textbf{Αυχενική δυσκαμψία} (δυσκολία στην κάμψη του αυχένα προς το στήθος).
\item Φωτοφοβία (ενόχληση στο φως), ναυτία, έμετοι, σύγχυση ή υπνηλία.
\item[] \textbf{Σημεία πίεσης αεραγωγού:}
\item \textbf{Δυσκολία στην αναπνοή} ή αίσθημα πνιγμού.
\item \textbf{Ταχέως αυξανόμενο οίδημα στον λαιμό} ή αίσθημα πίεσης/σφιξίματος.
\item \textbf{Αδυναμία κατάποσης} υγρών ή σιέλου.
\end{UrgentBox}

\textbf{Επικοινωνήστε άμεσα με τον ιατρό εάν εμφανίσετε:}
\begin{itemize}
\item Νέα ή επιδεινούμενη αδυναμία στα άνω ή κάτω άκρα.
\item Νέο μούδιασμα ή απώλεια αισθητικότητας στα άκρα.
\item Διαταραχή της βάδισης ή της ισορροπίας, ή αδεξιότητα στα χέρια (απώλεια λεπτών κινήσεων).
\item Ακράτεια ούρων/κοπράνων ή αδυναμία ούρησης.
\item Επιδεινούμενη ή επίμονη δυσκολία κατάποσης, ή επίμονο βράχνιασμα.
\item Έντονο ή επιδεινούμενο πόνο.
\item Ερυθρότητα, οίδημα ή εκροή από οποιοδήποτε τραύμα.
\item Πόνο, οίδημα ή ερυθρότητα σε κάτω άκρο (πιθανή φλεβοθρόμβωση).
% \item Επίμονη δυσκοιλιότητα παρά τη λήψη του καθαρτικού.
\end{itemize}

\newpage
\begin{center}
\LARGE\textbf{Οδηγίες Φαρμακευτικής Αγωγής}
\end{center}
\vspace{0.5cm}

\begin{MedTable}
\DrugEntry{inj Clexane}{4000IU}{1 x 1}{10}{αντιπηκτικές ενέσεις υποδορίως, την ίδια ώρα κάθε ημέρα}
\DrugEntry{tb Zinadol}{500mg}{1 x 2}{5}{αντιβίωση}
\DrugEntry{tb Apotel}{1000mg}{2 ή 3 την ημέρα}{}{παυσίπονο, μέγιστη δόση 3 δισκία το 24ωρο}
\DrugEntry{caps Omeprazole}{40mg}{1 x 2}{10}{προ φαγητού}
% ---- Drug library: uncomment the lines you need ---------------------
% -- Analgesics
% \DrugEntry{tb Paracetamol}{1000mg}{1 x 3}{10}{}
% \DrugEntry{tb Apotel}{1000mg}{2 ή 3 την ημέρα}{}{παυσίπονο, μέγιστη δόση 3 δισκία το 24ωρο}
% \DrugEntry{tb Depon}{500mg}{1 ή 2, έως 3 φορές την ημέρα}{}{εφόσον υπάρχει πόνος -- μέγιστο 3g το 24ωρο}
% \DrugEntry{tb Tramal}{50mg}{1 x 2 ή 3}{6}{}
% \DrugEntry{tb Zaldiar}{}{2 ή 3 την ημέρα}{}{εφόσον υπάρχει έντονος πόνος}
% -- NSAIDs
% \DrugEntry{tb Movatec}{7.5mg}{1 x 2}{10}{μετά το φαγητό}
% \DrugEntry{tb Movatec}{15mg}{1 x 1}{1}{μετά το φαγητό}
% \DrugEntry{tb Naprosyn}{500mg}{1 x 2}{6}{μετά το φαγητό}
% -- Steroids (tapers)
% \DrugEntry{tb Medrol}{4mg}{1 x 3\newline 1 x 2\newline 1 x 1\newline STOP}{}{για 5 ημέρες\newline για 5 ημέρες\newline για 5 ημέρες}
% \DrugEntry{tb Medrol}{4mg}{2 x 3\newline 1 x 3\newline 1 x 2\newline 1 x 1\newline STOP}{}{για 3 ημέρες\newline για 3 ημέρες\newline για 3 ημέρες\newline για 3 ημέρες}
% \DrugEntry{syr Dexaton}{2mg/5ml}{10cc x 3}{10}{μετά το φαγητό}
% \DrugEntry{tb Dexamethasone}{4mg}{1 x 3}{}{μετά το φαγητό}
% -- Antibiotics
% \DrugEntry{tb Zinadol}{500mg}{1 x 2}{5}{αντιβίωση}
% \DrugEntry{tb Klaricid}{500mg}{1 x 2}{7}{αντιβίωση}
% \DrugEntry{caps Amoxicillin}{500mg}{1 x 3}{7}{μετά το φαγητό}
% -- Gastroprotection
% \DrugEntry{caps Omeprazole}{40mg}{1 x 2}{10}{προ φαγητού}
% \DrugEntry{caps Omeprazole}{20mg}{1 x 2}{10}{15 λεπτά πριν το φαγητό}
% \DrugEntry{caps Mostrelan}{40mg}{1 x 1}{10}{προ φαγητού}
% -- Anticoagulants
% \DrugEntry{inj Clexane}{4000IU}{1 x 1}{10}{αντιπηκτικές ενέσεις υποδορίως, την ίδια ώρα κάθε ημέρα}
% \DrugEntry{inj Innohep}{0,45ml}{1 x 1}{10}{αντιπηκτικές ενέσεις υποδορίως, την ίδια ώρα κάθε ημέρα}
% \DrugEntry{inj Fraxiparine}{2850IU}{1 x 1}{10}{αντιπηκτικές ενέσεις υποδορίως, την ίδια ώρα κάθε ημέρα}
% -- Neuropathic pain, antiepileptics
% \DrugEntry{tb Lyrica}{50mg}{1 x 1 (βράδυ)}{4}{μετά την 4η ημέρα, 1 x 2}
% \DrugEntry{caps Lyrica}{150mg}{1 x 2}{}{πρωί -- βράδυ, με τροφή. Μην διακόπτεται απότομα}
% \DrugEntry{tb Neurontin}{400mg}{2 x 2}{}{}
% \DrugEntry{tb Tegretol}{200mg}{1 x 3}{}{με τροφή, σταθερές ώρες. \textbf{Μην διακόπτεται απότομα}}
% \DrugEntry{caps Cymbalta}{60mg}{1 x 2}{}{}
% \DrugEntry{tb Keppra}{1000mg}{1 x 3}{}{}
% \DrugEntry{tb Depakine}{400mg}{1 x 3}{}{θα γίνει προοδευτική μείωση σε συνεννόηση με τον θεράποντα ιατρό}
% -- Other
% \DrugEntry{sir Duphalac}{5ml}{1 ή 2 την ημέρα}{}{καθαρτικό, ανάλογα με τη λειτουργία του εντέρου -- στόχος: μαλακές κενώσεις χωρίς προσπάθεια}
% \DrugEntry{tb Xanax}{0,25mg}{1 x 2-3}{}{μέχρι 3 την ημέρα επί εκνευρισμού}
% \DrugEntry{tb Escitalopram}{10mg}{1 x 1}{}{}
% \DrugEntry{spray Vioplex-T ή Pulvo47}{}{1 x 2}{}{στο τραύμα, και να καλύπτεται μετά}
% ---------------------------------------------------------------------
\end{MedTable}

% Notes under the table (uncomment the ones you need)
\textbf{Σημείωση:} Μην λαμβάνετε αντιφλεγμονώδη φάρμακα (π.χ.\ Voltaren, Brufen, Ponstan, Naprosyn) χωρίς ιατρική οδηγία.
% \textbf{Σημείωση:} Μην διακόψετε και μην μειώσετε μόνος/μόνη σας τα Τιμή 8. Οι δόσεις θα μειωθούν σταδιακά κατόπιν ιατρικής οδηγίας.
% \textbf{Σημείωση:} Ενημερώστε κάθε ιατρό και φαρμακοποιό ότι λαμβάνετε Τιμή 9, πριν από τη χορήγηση οποιουδήποτε νέου φαρμάκου.
% \textbf{Σημείωση:} Η υπόλοιπη φαρμακευτική σας αγωγή (π.χ.\ αντιυπερτασική) συνεχίζεται \textbf{αμετάβλητη}.

\end{document}
